\section{Introduction}
\label{sec:introduction}

\IEEEPARstart{U}{rban} policy is enacted on parcels and districts, but its
consequences travel. When a municipality upzones a corridor, relaxes parking
minimums, or captures land value to fund transit, the resulting change in rents,
land values and accessibility does not stop at the boundary of the treated
zone. It diffuses to adjacent neighbourhoods through housing search, commuting
patterns, developer expectations and the local price of land. This diffusion is
not a nuisance to be controlled away; for many of the questions planners
actually ask---\emph{who is displaced, and where?}---it is the phenomenon of
interest.

This poses a problem for quantitative evaluation. The standard potential
outcomes framework rests on the Stable Unit Treatment Value Assumption
(SUTVA)~\cite{rubin1980}, which requires that a unit's outcome depend only on
its own treatment assignment. Spatial spillover violates SUTVA by construction.
A tract adjacent to an upzoned corridor is, in a meaningful sense, partially
treated even though it received no intervention. Estimators that ignore this
attribute the neighbour's rent increase either to noise or, worse, to the
control condition, biasing the estimated effect of the policy toward zero.

\subsection{Ex-Ante Comparison, Not Point Prediction}

A second and less-discussed problem concerns the framing of the modelling task
itself. Much of the machine learning literature on urban and mobility data
frames the objective as forecasting: predict next month's rents, next hour's
traffic, next year's population. Spatio-temporal graph networks have been
notably successful at exactly this~\cite{li2018dcrnn,yu2018stgcn}. But a
planning agency deciding between an inclusionary housing mandate and a
land-value capture scheme is not asking what rents will be in 2031. It is
asking which of several counterfactual worlds is preferable, and for whom.
These are different questions with different evidentiary standards. A point
forecast can be accurate and still be useless for a policy choice if it cannot
be conditioned on an intervention that has not yet happened.

We therefore adopt an explicitly \emph{ex-ante comparative} scope rule. The
quantity of interest is not $\hat{Y}_{i,t+1}$ in absolute terms but the
contrast between $\hat{Y}_{i,t+1}$ under an intervention and
$\hat{Y}_{i,t+1}$ under the untreated counterfactual, together with the
corresponding contrast at each of unit $i$'s neighbours. Forecast accuracy
enters our evaluation not as the goal but as a necessary condition: a model
whose conditional expectations are badly calibrated cannot be trusted to
produce meaningful contrasts between counterfactuals.

\subsection{Approach}

We introduce the Causal Graph Counterfactual Exposure Network (\CGCEN{}), which
makes the interference structure explicit in the architecture rather than
implicit in a learned representation. Following the partial interference
formulation of Hudgens and Halloran~\cite{hudgens2008}, we assume a unit's
outcome depends on its own treatment and on a summary of its neighbours'
treatments mediated by the graph. \CGCEN{} factorises the one-step-ahead
prediction as
\begin{equation}
  \hat{Y}_{i,t+1} = \underbrace{\Phi(X_{i,t}, P_{i,t})}_{\text{direct}}
  + \underbrace{\sum_{j \in \mathcal{N}(i)} W_{ij} \cdot
  \Psi(X_{j,t}, P_{j,t})}_{\text{spatial interference}} ,
  \label{eq:decomposition-intro}
\end{equation}
where $\Phi$ estimates the Individual Treatment Effect (ITE) and $\Psi$,
aggregated over neighbours with row-normalised weights $W_{ij}$, captures the
Spatial Interference Effect (STE). The separation is enforced structurally:
$\Phi$ consumes a node-local temporal encoding that never passes through a
message-passing operator, so by construction it cannot launder neighbour
information into the direct term. Both effects are recovered by contrasting
each branch against its own $P = 0$ counterfactual, which makes the ITE exactly
zero on untreated units and the STE exactly zero on an edgeless graph---two
identities we verify numerically as regression tests.

Because credible spatial policy data with known ground truth is scarce, we
evaluate on a transparent synthetic city whose data-generating process we
control. This is a deliberate methodological choice. It lets us ask a question
that observational data cannot answer: \emph{when the spillover is known to
exist, does the graph model actually find it?} The answer turns out to be
conditional, and that condition is our main contribution.

\subsection{Related Work}

Our work sits at the intersection of three literatures.

\emph{Causal inference under interference} relaxes SUTVA by positing that a
unit's outcome depends on others' assignments. Hudgens and
Halloran~\cite{hudgens2008} formalise partial interference within disjoint
clusters; Ogburn and VanderWeele~\cite{ogburn2014} give a graphical account;
and Forastiere et al.~\cite{forastiere2021} extend identification to
observational network settings with an explicit exposure mapping. We adopt
their exposure-mapping formulation but embed it in a learned spatio-temporal
model rather than a parametric outcome regression.

\emph{Neural counterfactual estimation} learns representations for treatment
effect estimation, typically under SUTVA. Shalit et al.~\cite{shalit2017cfr}
give generalisation bounds for individual treatment effects, and Wager and
Athey~\cite{wager2018causalforest} provide asymptotic inference for forest-based
heterogeneous effects. Ma and Tresp~\cite{ma2021netdeconf} extend neural
counterfactual estimation to networked interference, which is closest to our
setting; our contribution relative to that line is the architectural
identification of the direct/interference split and the explicit temporal
dimension.

\emph{Spatio-temporal graph learning} has been driven largely by traffic
forecasting. Kipf and Welling~\cite{kipf2017gcn} introduced the graph
convolution we take as a reference point, and DCRNN~\cite{li2018dcrnn} and
STGCN~\cite{yu2018stgcn} established the recurrent and convolutional
spatio-temporal templates. These are point-forecasting models evaluated on
observational series; neither exposes a treatment channel or a counterfactual
contrast, which is what the policy-comparison setting requires.

On the measurement side, the combined housing-and-transportation affordability
framing follows the H$+$T Index~\cite{cnt2023ht}, and our displacement
indicators are informed by the empirical displacement literature, particularly
Zuk et al.~\cite{zuk2018gentrification}, who document how public investment can
simultaneously raise accessibility and displacement risk---the tension our
$\HT$ index is designed to make visible.

\subsection{Contributions}

\begin{enumerate}
  \item \textbf{An architecturally identified interference split.} We give a
  spatio-temporal architecture in which the direct and interference channels
  are separated by construction rather than by regularisation, together with
  counterfactual contrasts that recover ITE and STE with exact degenerate-case
  behaviour (Section~\ref{sec:methodology}).

  \item \textbf{A level-aware displacement index.} We show that a
  growth-rate-only Displacement Pressure Index is structurally incapable of
  registering a sustained policy effect: a permanent level shift produces a
  single transient spike and then returns to zero. We add a cumulative
  rent-burden term that restores sensitivity to persistent burden, increasing
  measured policy contrasts by more than two orders of magnitude
  (Sections~\ref{sec:methodology} and~\ref{sec:experiments}).

  \item \textbf{A conditional result on when graphs help.} \CGCEN{} improves
  significantly over a linear spatial lag model under active policy
  intervention pulses ($p = 0.0015$--$0.0033$), but \emph{not} in the untreated
  equilibrium ($p = 0.127$). We report this boundary explicitly rather than
  pooling it away, because the pooled comparison would overstate the
  method's generality.

  \item \textbf{A negative result with a mechanism.} Under a generator in which
  spillover arrives as a single permanent shock, no model we tested---including
  ours, and including the spatial lag model whose learned spatial coefficient
  converged to $\rho \approx 0$---derived measurable benefit from the graph.
  Removing edges \emph{improved} accuracy. We diagnose the cause: permanent
  early spillover is absorbed into each node's own level and becomes a static
  intercept, indistinguishable from node-specific effects. We show that making
  the shock process stochastic and recurrent restores identifiability
  (Section~\ref{sec:data-engine}).

  \item \textbf{A reproducible benchmark.} We release the generator, a
  four-model baseline suite, a ten-seed harness with confidence intervals and
  paired significance tests, and topology ablations (real, edgeless, randomly
  rewired), together with all figure and table assets used in this manuscript.
\end{enumerate}

The fourth contribution deserves emphasis. It is the kind of finding that is
usually discarded during model development, and discarding it here would have
produced a paper claiming that a graph neural network captures urban spillover
when in fact the benchmark contained no dynamic spillover for it to capture.
Reporting the boundary between the regimes is, we argue, more useful to the
field than reporting only the regime in which the method wins.

\subsection{Organisation}

Section~\ref{sec:methodology} formalises the graph, the causal decomposition
and the two policy indices. Section~\ref{sec:data-engine} describes the
synthetic data engine, the dynamic shock process and the logistic stabilisation
that keeps long-horizon trajectories bounded.
Section~\ref{sec:experiments} reports the ten-seed benchmark, significance
tests and topology ablations. Section~\ref{sec:discussion} discusses
limitations---chiefly the absence of a full-scale empirical
evaluation---and outlines the path to external validation.
